%% kapa: exam=2 round=25 year=2014 period=1 subject=appraisal_practice subject_ko=감정평가실무 seq=2 number=2 points=30
%% source: Q-Net 감정평가사 기출문제 (cst00309) artlSeq=5109209 file=extracted/1008/5109209/2014년 제25회 감정평가사 2차시험 기출문제/1교시(감정평가실무).hwp
%% conversion: kordoc 4.24.51 -> tools/gen2.py (auto); review=proofread-sonnet; flags=none
\begin{longquestion}{【문제 2】}
서울특별시 A구가 도시계획시설 도로개설사업으로 소유자 甲의 토지를 협의 취득하였으나 이후 당해 토지가 B공사가 시행하는 H택지개발사업에 편입되어 환매권이 발생하였다. 그러나 사업시행자가 원소유자 甲에게 환매권이 발생한 사실을 통지나 공고를 하지 아니하여 결국 환매권이 상실되었다. 이에 원소유자 甲은 환매권 발생 통지의무 해태로 인한 손해배상소송을 제기하였다. 원소유자 甲의 환매권 상실로 인한 손해액을 다음 물음에 따라 구하시오. (30점) \par
\medskip\noindent 물음 1) 환매권 상실 당시의 토지 평가에 적용할 비교표준지 기호 및 적용공시지가(연도)와 그 선정이유 (5점) \par
\medskip\noindent 물음 2) 환매권 상실 당시의 토지 평가금액 (5점) \par
\medskip\noindent 물음 3) 환매권을 행사하였을 경우 반환하여야 할 환매금액 (15점) \par
\medskip\noindent 물음 4) 환매권 상실로 인한 손해액 (5점) \par
\textbf{<자료 1> 개 요} \par
1. 지급한 보상금액은 토지 56,700,000원, 지장물 5,400,000원이며, 환매 토지의 소유권이전일은 2001. 9. 20임 \par
2. H택지개발사업의 사업인정고시 의제일은 2007. 10. 27임 \par
\textbf{<자료 2> 토지의 개황} \par
1. 소재지 및 면적 : 서울특별시 A구 B동 119번지, 315m\textsuperscript{2} \par
2. 도로개설사업 편입 당시 : 전, 부정형, 완경사, 맹지 \par
3. 가격조사완료일 : 2014. 9. 20 \par
4. 환매대상토지는 2004년 5월 중에 용도지역이 자연녹지지역에서 제2종일반주거지역으로 변경되었음 \par
\textbf{<자료 3> 표준지공시지가} \par
1. 인근 표준지공시지가의 특성항목 \par
{\small\centering
\begin{tabular}{|c|c|c|c|c|c|c|c|}
\hline
\makecell{기호} & \makecell{소재지} & \makecell{지목} & \makecell{면적(m\textsuperscript{2})} & \makecell{이용상황} & \makecell{용도지역} & \makecell{도로교통} & \makecell{형상\newline{}지세} \\
\hline
\makecell{1} & \makecell{A구 B동\newline{}16} & \makecell{전} & \makecell{865} & \makecell{전} & \makecell{자연녹지} & \makecell{맹지} & \makecell{부정형\newline{}완경사} \\
\hline
\makecell{2} & \makecell{A구 B동\newline{}255-1} & \makecell{대} & \makecell{540} & \makecell{단독주택} & \makecell{자연녹지} & \makecell{세로(가)} & \makecell{부정형\newline{}평지} \\
\hline
\makecell{3} & \makecell{A구 B동\newline{}306} & \makecell{전} & \makecell{306} & \makecell{전} & \makecell{자연녹지} & \makecell{세로(가)} & \makecell{부정형\newline{}완경사} \\
\hline
\makecell{4} & \makecell{A구 B동\newline{}381-5} & \makecell{전} & \makecell{413} & \makecell{주거나지} & \makecell{2종일주} & \makecell{세로(불)} & \makecell{부정형\newline{}완경사} \\
\hline
\makecell{5} & \makecell{A구 B동\newline{}651} & \makecell{대} & \makecell{248} & \makecell{단독주택\newline{}(다가구)} & \makecell{2종일주} & \makecell{세로(가)} & \makecell{가장형\newline{}평지} \\
\hline
\makecell{6} & \makecell{A구 C동\newline{}381-5} & \makecell{전} & \makecell{243} & \makecell{단독주택} & \makecell{2종일주} & \makecell{세로(가)} & \makecell{사다리\newline{}평지} \\
\hline
\end{tabular}\par}

1) 기호2, 4, 5 표준지공시지가는 도로개설사업에 따른 가격변동이 있는 것으로 조사되었음 \par
2) 기호3, 5 표준지공시지가는 2004년 신규표준지임 \par
3) 기호6 표준지공시지가는 환매 대상 토지와 용도지역, 지목 및 이용상황의 변경과정이 유사함 \par
2. 인근 표준지 공시지가의 연도별 공시가격 \par
(단위 : 원/m\textsuperscript{2}) \par
{\small\centering
\begin{tabular}{|c|c|c|c|c|c|c|}
\hline
\makecell{기호} & \makecell{2001년} & \makecell{2002년} & \makecell{2007년} & \makecell{2008년} & \makecell{2011년} & \makecell{2012년} \\
\hline
\makecell{1} & \makecell{100,000} & \makecell{106,000} & \makecell{137,000} & \makecell{129,000} & \makecell{145,000} & \makecell{158,000} \\
\hline
\makecell{2} & \makecell{200,000} & \makecell{210,000} & \makecell{283,000} & \makecell{270,000} & \makecell{300,000} & \makecell{330,000} \\
\hline
\makecell{3} & \makecell{-} & \makecell{-} & \makecell{170,000} & \makecell{165,000} & \makecell{185,000} & \makecell{202,000} \\
\hline
\makecell{4} & \makecell{270,000} & \makecell{284,000} & \makecell{380,000} & \makecell{370,000} & \makecell{410,000} & \makecell{450,000} \\
\hline
\makecell{5} & \makecell{-} & \makecell{-} & \makecell{530,000} & \makecell{510,000} & \makecell{570,000} & \makecell{620,000} \\
\hline
\makecell{6} & \makecell{360,000} & \makecell{380,000} & \makecell{500,000} & \makecell{490,000} & \makecell{550,000} & \makecell{600,000} \\
\hline
\end{tabular}\par}

\textbf{<자료 4> 지가변동률(서울특별시 A구)} \par
{\small\centering
\begin{tabular}{|c|c|c|}
\hline
\multirow{2}{*}{\makecell{기간}} & \multicolumn{2}{|c|}{\makecell{변동률(\%)}} \\
\cline{2-3}
 & \makecell{주거지역} & \makecell{녹지지역} \\
\hline
\makecell{2001. 1. 1 ～ 2001. 9. 20} & \makecell{0.100} & \makecell{0.200} \\
\hline
\makecell{2001. 1. 1 ～ 2007. 10. 27} & \makecell{0.300} & \makecell{0.400} \\
\hline
\makecell{2001. 1. 1 ～ 2011. 9. 20} & \makecell{3.500} & \makecell{4.500} \\
\hline
\makecell{2007. 1. 1 ～ 2007. 10. 27} & \makecell{- 0.001} & \makecell{- 0.002} \\
\hline
\makecell{2007. 1. 1 ～ 2011. 9. 20} & \makecell{5.000} & \makecell{7.000} \\
\hline
\makecell{2011. 1. 1 ～ 2011. 9. 20} & \makecell{0.060} & \makecell{0.090} \\
\hline
\end{tabular}\par}

\textbf{<자료 5> 지역요인 및 개별요인 비교} \par
1. 본건 토지와 표준지공시지가 기호1 ～ 기호6은 인근지역에 위치하므로 지역요인은 대등함 \par
2. 개별요인 비교치 \par
{\small\centering
\begin{tabular}{|c|c|c|c|c|c|c|}
\hline
\makecell{구분} & \makecell{표준지1} & \makecell{표준지2} & \makecell{표준지3} & \makecell{표준지4} & \makecell{표준지5} & \makecell{표준지6} \\
\hline
\makecell{개별요인 비교치} & \makecell{1.010} & \makecell{1.020} & \makecell{1.030} & \makecell{1.040} & \makecell{1.050} & \makecell{1.060} \\
\hline
\end{tabular}\par}

\textbf{<자료 6> 그 밖의 요인 보정치} \par
본건 평가에 적용할 그 밖의 요인 보정치는 표준지공시지가 기호1 ∼ 기호6 공히 1.30으로 산정되었음 \par
\textbf{<자료 7> 기타} \par
1. 환매대상 토지는 환매권 상실 당시 H택지개발지구에 편입되어 조성공사 중인 바 지적확인이 곤란한 상태이며, 인근의 표준적인 이용상황은 주거용(다가구주택)임 \par
2. 상기의 B공사는「공공기관의 운영에 관한 법률」제5조 제3항 제1호의 공공기관임 \par
3. 상기의 H택지개발사업은「공익사업을 위한 토지 등의 취득 및 보상에 관한 법률」제4조 제5호에 규정된 공익사업임 \par
\end{longquestion}
